\documentclass[12pt,a4paper]{article}

% ─── Pacotes ──────────────────────────────────────────────────────────────────
\usepackage[utf8]{inputenc}
\usepackage[T1]{fontenc}
\usepackage[portuguese]{babel}
\usepackage{geometry}
\geometry{margin=2.5cm}
\usepackage{graphicx}
\usepackage{hyperref}
\hypersetup{colorlinks=true, linkcolor=blue, urlcolor=blue, citecolor=blue}
\usepackage{booktabs}
\usepackage{longtable}
\usepackage{array}
\usepackage{xcolor}
\usepackage{listings}
\usepackage{enumitem}
\usepackage{fancyhdr}
\usepackage{titlesec}
\usepackage{parskip}
\usepackage{mdframed}
\usepackage{tikz}
\usetikzlibrary{arrows.meta, shapes.rounded, positioning}

% ─── Cores ────────────────────────────────────────────────────────────────────
\definecolor{primaryblue}{RGB}{13,110,253}
\definecolor{lightgray}{RGB}{248,249,250}
\definecolor{codegray}{RGB}{240,240,240}
\definecolor{successgreen}{RGB}{25,135,84}
\definecolor{warningyellow}{RGB}{255,193,7}
\definecolor{dangerred}{RGB}{220,53,69}

% ─── Estilo de código ─────────────────────────────────────────────────────────
\lstdefinestyle{bashstyle}{
  backgroundcolor=\color{codegray},
  basicstyle=\ttfamily\small,
  breaklines=true,
  frame=single,
  rulecolor=\color{gray!40},
  xleftmargin=8pt,
  xrightmargin=8pt,
  aboveskip=6pt,
  belowskip=6pt,
}
\lstset{style=bashstyle}

% ─── Caixas de destaque ───────────────────────────────────────────────────────
\newmdenv[
  backgroundcolor=lightgray,
  linecolor=primaryblue,
  linewidth=2pt,
  topline=false, rightline=false, bottomline=false,
  leftmargin=0pt, rightmargin=0pt,
  innerleftmargin=10pt, innerrightmargin=10pt,
  innertopmargin=6pt, innerbottommargin=6pt,
]{infobox}

% ─── Cabeçalho ────────────────────────────────────────────────────────────────
\pagestyle{fancy}
\fancyhf{}
\rhead{\textcolor{gray}{PAW --- Grupo 13}}
\lhead{\textcolor{gray}{Marketplace de Supermercados}}
\rfoot{\thepage}
\renewcommand{\headrulewidth}{0.4pt}

% ─── Secções coloridas ────────────────────────────────────────────────────────
\titleformat{\section}{\Large\bfseries\color{primaryblue}}{}{0em}{\thesection\quad}[\titlerule]
\titleformat{\subsection}{\large\bfseries}{}{0em}{\thesubsection\quad}
\titleformat{\subsubsection}{\normalsize\bfseries}{}{0em}{\thesubsubsection\quad}

% ─── Título ───────────────────────────────────────────────────────────────────
\title{
  \vspace{-1cm}
  \Large \textbf{Programação Avançada para a Web}\\[0.4cm]
  \large Trabalho Prático --- Grupo 13 \quad|\quad ESTG --- P.Porto\\[0.2cm]
  \rule{\linewidth}{1pt}\\[0.4cm]
  \LARGE \textbf{Marketplace de Supermercados Locais}\\[0.3cm]
  \large Relatório de Progresso --- Milestone 1\\[0.2cm]
  \rule{\linewidth}{1pt}
}
\author{Grupo 13}
\date{Abril 2026}

% ══════════════════════════════════════════════════════════════════════════════
\begin{document}

\maketitle
\thispagestyle{empty}
\newpage

\tableofcontents
\newpage

% ══════════════════════════════════════════════════════════════════════════════
\section{Introdução}

O presente relatório descreve o trabalho desenvolvido pelo Grupo 13 no âmbito da unidade curricular de \textbf{Programação Avançada para a Web (PAW)}. O projeto consiste no desenvolvimento de uma plataforma \textit{marketplace} para supermercados locais, que centraliza a gestão de lojas, produtos, encomendas, estafetas e clientes através de uma aplicação web completa.

\subsection{Objetivos da Plataforma}

A plataforma oferece um conjunto alargado de funcionalidades, organizadas por perfil de utilizador:

\begin{itemize}[leftmargin=*, label=\textcolor{primaryblue}{\textbullet}]
  \item Registo e aprovação de supermercados com notificação por email;
  \item Gestão de produtos por categoria, com suporte a imagens e controlo de stock;
  \item Registo de vendas presenciais com carrinho, cupões de desconto e método de entrega;
  \item Fluxo completo de encomendas com estados e notificações automáticas por email;
  \item Rede de estafetas para aceitação e conclusão de entregas ao domicílio;
  \item Sistema de cupões de desconto por percentagem, com validade e limite de usos;
  \item Mapa interativo com geolocalização automática dos supermercados;
  \item API REST completamente documentada com Swagger/OpenAPI 3.0.
\end{itemize}

\subsection{Perfis de Utilizador}

A plataforma define quatro perfis com permissões diferenciadas e portais independentes:

\begin{center}
\begin{tabular}{>{\ttfamily}lp{9cm}}
\toprule
\textbf{Perfil} & \textbf{Responsabilidades} \\
\midrule
admin       & Gestão global: aprovação de lojas, gestão de utilizadores, categorias, produtos, encomendas, cupões e mapa \\
supermarket & Gestão da sua loja: produtos, vendas presenciais, encomendas recebidas, cupões e definições \\
courier     & Aceitação de entregas disponíveis e marcação como entregue \\
customer    & Cliente final — registo e histórico de encomendas (M2) \\
\bottomrule
\end{tabular}
\end{center}

% ══════════════════════════════════════════════════════════════════════════════
\section{Arquitetura e Tecnologias}

\subsection{Stack Tecnológica}

\begin{center}
\begin{tabular}{ll}
\toprule
\textbf{Componente}      & \textbf{Tecnologia / Versão} \\
\midrule
Runtime                  & Node.js $\geq$ 18 \\
Framework web            & Express 4 \\
Base de dados            & MongoDB + Mongoose 9 \\
Motor de templates       & EJS \\
Autenticação             & JWT em cookie \texttt{HttpOnly} (jsonwebtoken 9) \\
Upload de ficheiros      & Multer 2 \\
Email                    & Nodemailer 8 (Gmail SMTP) \\
Segurança                & express-mongo-sanitize \\
Geocodificação           & Nominatim (OpenStreetMap API) \\
Mapa interativo          & Leaflet.js + OpenStreetMap \\
Documentação API         & Swagger UI + swagger-jsdoc (OpenAPI 3.0) \\
UI / CSS                 & Bootstrap 5.3 + Bootstrap Icons 1.11 \\
Alertas frontend         & SweetAlert2 \\
\bottomrule
\end{tabular}
\end{center}

\subsection{Arquitetura MVC}

O projeto segue uma arquitetura \textbf{MVC (Model-View-Controller)} com separação clara de responsabilidades:

\begin{itemize}[leftmargin=*, label=\textcolor{primaryblue}{\textbullet}]
  \item \textbf{Models} (\texttt{/models}): esquemas Mongoose para \texttt{User}, \texttt{Supermarket}, \texttt{Product}, \texttt{Category}, \texttt{Order} e \texttt{Coupon};
  \item \textbf{Views} (\texttt{/views}): templates EJS organizados por perfil (\texttt{admin/}, \texttt{supermarket/}, \texttt{courier/}, \texttt{auth/});
  \item \textbf{Controllers} (\texttt{/controllers}): lógica de negócio separada em controllers de API e controllers de vistas;
  \item \textbf{Routes} (\texttt{/routes}): mapeamento de rotas por domínio --- \texttt{api.js}, \texttt{admin.js}, \texttt{supermarket.js}, \texttt{courier.js};
  \item \textbf{Middlewares} (\texttt{/middlewares}): \texttt{verifyToken} (validação JWT), \texttt{verifyRole} (controlo de acesso por perfil) e \texttt{verifySupermarketStatus} (bloqueia supermercados não aprovados);
  \item \textbf{Utils} (\texttt{/utils}): serviços reutilizáveis para email, geocodificação, Swagger e acesso à base de dados.
\end{itemize}

\subsection{Segurança}

\begin{itemize}[leftmargin=*, label=\textcolor{primaryblue}{\textbullet}]
  \item Tokens JWT armazenados em cookies \texttt{HttpOnly} — inacessíveis via JavaScript no browser;
  \item Passwords com hash \texttt{bcrypt} (salt rounds: 10);
  \item \texttt{express-mongo-sanitize} para prevenção de injeção NoSQL;
  \item Controlo de acesso por role em todas as rotas protegidas;
  \item Validação de propriedade em operações sensíveis (ex: cancelamento de encomenda);
  \item \textbf{Conta superadmin} protegida: criada automaticamente na primeira inicialização via \texttt{utils/seed.js}, com flag \texttt{isSuperAdmin} na base de dados — invisível na tabela de admins, imune a edição e eliminação por qualquer utilizador.
\end{itemize}

% ══════════════════════════════════════════════════════════════════════════════
\section{Funcionalidades Implementadas}

\subsection{Rotas de Interface (EJS)}

\subsubsection{Autenticação --- \texttt{/auth}}

\begin{center}
\begin{tabular}{lll}
\toprule
\textbf{Método} & \textbf{Rota} & \textbf{Descrição} \\
\midrule
GET  & \texttt{/auth/login}    & Página de login \\
POST & \texttt{/auth/login}    & Submissão de credenciais \\
GET  & \texttt{/auth/logout}   & Logout (elimina cookie JWT) \\
GET  & \texttt{/auth/register} & Página de registo de supermercado \\
POST & \texttt{/auth/register} & Submissão de registo \\
\bottomrule
\end{tabular}
\end{center}

\subsubsection{Painel de Administração --- \texttt{/admin}}

Todas as rotas requerem perfil \texttt{admin}.

\begin{center}
\begin{tabular}{ll}
\toprule
\textbf{Rota} & \textbf{Descrição} \\
\midrule
\texttt{/admin/dashboard}           & Dashboard com estatísticas globais (lojas, utilizadores, encomendas) \\
\texttt{/admin/approvals}           & Aprovação e rejeição de supermercados pendentes \\
\texttt{/admin/users}               & Gestão de todos os utilizadores (tabs por perfil) \\
\texttt{/admin/categories}          & Criação e gestão de categorias de produtos \\
\texttt{/admin/products}            & Visão global de todos os produtos \\
\texttt{/admin/product/:id}         & Detalhe e edição de produto \\
\texttt{/admin/orders}              & Listagem de todas as encomendas \\
\texttt{/admin/order/:id}           & Detalhe de encomenda com histórico \\
\texttt{/admin/supermarket/:id}     & Detalhe de supermercado com estatísticas \\
\texttt{/admin/user/:id}            & Detalhe de utilizador \\
\texttt{/admin/map}                 & Mapa interativo com supermercados georreferenciados \\
\texttt{/admin/coupons}             & Listagem de todos os cupões da plataforma \\
\texttt{/admin/coupon/:id}          & Detalhe e edição de cupão \\
\bottomrule
\end{tabular}
\end{center}

\subsubsection{Portal do Supermercado --- \texttt{/supermarket}}

Todas as rotas requerem perfil \texttt{supermarket} e estado \texttt{approved}.

\begin{center}
\begin{tabular}{ll}
\toprule
\textbf{Rota} & \textbf{Descrição} \\
\midrule
\texttt{/supermarket/dashboard}     & Dashboard com métricas e encomendas recentes \\
\texttt{/supermarket/settings}      & Definições da loja (nome, morada, horário, entregas) e conta \\
\texttt{/supermarket/products}      & Gestão de produtos (criar, editar, ativar/desativar) \\
\texttt{/supermarket/product/:id}   & Detalhe e edição de produto \\
\texttt{/supermarket/orders}        & Listagem e gestão de encomendas recebidas \\
\texttt{/supermarket/order/:id}     & Detalhe de encomenda e progressão de estado \\
\texttt{/supermarket/newsale}       & Registar nova venda presencial com carrinho \\
\texttt{/supermarket/coupons}       & Criação e gestão de cupões da loja \\
\bottomrule
\end{tabular}
\end{center}

\subsubsection{Portal do Estafeta --- \texttt{/courier}}

Todas as rotas requerem perfil \texttt{courier}.

\begin{center}
\begin{tabular}{ll}
\toprule
\textbf{Rota} & \textbf{Descrição} \\
\midrule
\texttt{/courier/dashboard}   & Dashboard com entrega ativa (se existir) \\
\texttt{/courier/deliveries}  & Lista de entregas disponíveis para aceitar \\
\texttt{/courier/history}     & Histórico de entregas concluídas \\
\texttt{/courier/order/:id}   & Detalhe da entrega e ação de conclusão \\
\bottomrule
\end{tabular}
\end{center}

\subsection{API REST --- \texttt{/api}}

A API devolve sempre JSON. A autenticação é feita via cookie \texttt{token} (JWT HttpOnly). A documentação interativa completa está disponível em \texttt{/api-docs} (Swagger UI).

\subsubsection{Utilizadores}

\begin{center}
\begin{tabular}{llll}
\toprule
\textbf{M} & \textbf{Rota} & \textbf{Acesso} & \textbf{Descrição} \\
\midrule
POST   & \texttt{/api/users/createAdmin}       & admin & Criar admin \\
GET    & \texttt{/api/users/listAdmins}        & admin & Listar admins \\
GET    & \texttt{/api/users/admin/:id}         & admin & Detalhe de admin \\
PUT    & \texttt{/api/users/updateAdmin/:id}   & admin & Editar admin \\
DELETE & \texttt{/api/users/deleteAdmin/:id}   & admin & Eliminar admin \\
\midrule
POST   & \texttt{/api/users/createCourier}     & admin & Criar estafeta \\
GET    & \texttt{/api/users/listCouriers}      & admin & Listar estafetas \\
GET    & \texttt{/api/users/courier/:id}       & admin & Detalhe de estafeta \\
PUT    & \texttt{/api/users/updateCourier/:id} & admin & Editar estafeta \\
DELETE & \texttt{/api/users/deleteCourier/:id} & admin & Eliminar estafeta \\
\midrule
POST   & \texttt{/api/customers/create}        & admin, supermarket & Criar cliente \\
GET    & \texttt{/api/customers/list}          & admin & Listar clientes \\
GET    & \texttt{/api/customers/:id}           & admin & Detalhe de cliente \\
GET    & \texttt{/api/customers/email/:email}  & admin, supermarket & Cliente por email \\
PUT    & \texttt{/api/customers/update/:id}    & admin & Editar cliente \\
DELETE & \texttt{/api/customers/delete/:id}    & admin & Eliminar cliente \\
\bottomrule
\end{tabular}
\end{center}

\subsubsection{Supermercados}

\begin{center}
\begin{tabular}{llll}
\toprule
\textbf{M} & \textbf{Rota} & \textbf{Acesso} & \textbf{Descrição} \\
\midrule
POST   & \texttt{/api/supermarkets/create}         & admin & Criar supermercado \\
GET    & \texttt{/api/supermarkets/list}           & admin & Listar todos \\
GET    & \texttt{/api/supermarkets/listPending}    & admin & Listar pendentes \\
GET    & \texttt{/api/supermarkets/map}            & público & Dados para o mapa \\
GET    & \texttt{/api/supermarkets/:id}            & admin & Detalhe \\
PUT    & \texttt{/api/supermarkets/approve/:id}    & admin & Aprovar \\
PUT    & \texttt{/api/supermarkets/reject/:id}     & admin & Rejeitar \\
PUT    & \texttt{/api/supermarkets/update/:id}     & admin & Editar (admin) \\
PUT    & \texttt{/api/supermarkets/me}             & supermarket & Editar a própria loja \\
DELETE & \texttt{/api/supermarkets/delete/:id}     & admin & Eliminar \\
\bottomrule
\end{tabular}
\end{center}

\subsubsection{Categorias e Produtos}

\begin{center}
\begin{tabular}{llll}
\toprule
\textbf{M} & \textbf{Rota} & \textbf{Acesso} & \textbf{Descrição} \\
\midrule
POST   & \texttt{/api/categories/create}       & admin & Criar categoria \\
GET    & \texttt{/api/categories/list}         & admin, supermarket & Listar \\
GET    & \texttt{/api/categories/:id}          & admin & Detalhe \\
PUT    & \texttt{/api/categories/update/:id}   & admin & Editar \\
DELETE & \texttt{/api/categories/delete/:id}   & admin & Eliminar \\
\midrule
POST   & \texttt{/api/product/create}          & admin, supermarket & Criar produto \\
GET    & \texttt{/api/product/list/me}         & supermarket & Produtos da minha loja \\
GET    & \texttt{/api/product/list/:id}        & admin, supermarket & Produtos por loja \\
GET    & \texttt{/api/product/:id}             & admin, supermarket & Detalhe \\
PUT    & \texttt{/api/product/update/:id}      & admin, supermarket & Editar \\
PUT    & \texttt{/api/product/toggle/:id}      & admin, supermarket & Ativar/desativar \\
DELETE & \texttt{/api/product/delete/:id}      & admin, supermarket & Eliminar \\
\bottomrule
\end{tabular}
\end{center}

\subsubsection{Encomendas}

\begin{center}
\begin{tabular}{llll}
\toprule
\textbf{M} & \textbf{Rota} & \textbf{Acesso} & \textbf{Descrição} \\
\midrule
POST   & \texttt{/api/orders/sale}                  & supermarket & Venda presencial \\
GET    & \texttt{/api/orders/list}                  & admin & Todas as encomendas \\
GET    & \texttt{/api/orders/supermarket/:id}       & admin, supermarket & Por supermercado \\
GET    & \texttt{/api/orders/customer/:id}          & admin, supermarket & Por cliente \\
GET    & \texttt{/api/orders/:id}                   & admin, supermarket & Detalhe \\
PUT    & \texttt{/api/orders/status/:id}            & admin, supermarket & Atualizar estado \\
PUT    & \texttt{/api/orders/cancel/:id}            & admin, supermarket & Cancelar \\
PUT    & \texttt{/api/orders/courier/accept/:id}    & courier & Aceitar entrega \\
PUT    & \texttt{/api/orders/courier/complete/:id}  & courier & Concluir entrega \\
\bottomrule
\end{tabular}
\end{center}

\subsubsection{Cupões}

\begin{center}
\begin{tabular}{llll}
\toprule
\textbf{M} & \textbf{Rota} & \textbf{Acesso} & \textbf{Descrição} \\
\midrule
POST   & \texttt{/api/coupons/create}          & supermarket & Criar cupão \\
GET    & \texttt{/api/coupons/list}            & admin, supermarket & Listar \\
GET    & \texttt{/api/coupons/:id}             & admin, supermarket & Detalhe \\
GET    & \texttt{/api/coupons/validate}        & supermarket & Validar cupão \\
PUT    & \texttt{/api/coupons/update/:id}      & admin, supermarket & Editar \\
PUT    & \texttt{/api/coupons/toggle/:id}      & admin, supermarket & Ativar/desativar \\
DELETE & \texttt{/api/coupons/delete/:id}      & admin, supermarket & Eliminar \\
\bottomrule
\end{tabular}
\end{center}

\subsection{Fluxo de Estados de Encomenda}

O ciclo de vida de uma encomenda passa pelos seguintes estados:

\vspace{0.5cm}
\begin{center}
\begin{tikzpicture}[
  node distance=1.2cm and 2cm,
  every node/.style={font=\small},
  state/.style={rounded rectangle, draw=primaryblue, fill=lightgray, minimum width=2.8cm, minimum height=0.7cm, align=center},
  arr/.style={-{Stealth}, thick, color=primaryblue},
  cancel/.style={-{Stealth}, thick, color=dangerred, dashed},
]
  \node[state] (pending)   {\texttt{pending}};
  \node[state, right=of pending] (confirmed) {\texttt{confirmed}};
  \node[state, right=of confirmed] (preparing) {\texttt{preparing}};
  \node[state, below right=0.8cm and 0.4cm of preparing] (awc) {\texttt{awaiting\_courier}};
  \node[state, right=1.2cm of awc] (delivering2) {\texttt{delivering}};
  \node[state, above right=0.8cm and 0.4cm of preparing] (delivering1) {\texttt{delivering}};
  \node[state, right=1.2cm of delivering1] (delivered1) {\texttt{delivered}};
  \node[state, right=1.2cm of delivering2] (delivered2) {\texttt{delivered}};

  \draw[arr] (pending) -- (confirmed);
  \draw[arr] (confirmed) -- (preparing);
  \draw[arr] (preparing) -- node[above, sloped, font=\tiny]{estafeta} (awc);
  \draw[arr] (preparing) -- node[above, sloped, font=\tiny]{levantamento} (delivering1);
  \draw[arr] (awc) -- (delivering2);
  \draw[arr] (delivering1) -- (delivered1);
  \draw[arr] (delivering2) -- (delivered2);

  \node[state, fill=dangerred!15, draw=dangerred, below=1.5cm of confirmed] (cancelled) {\texttt{cancelled}};
  \draw[cancel] (pending.south) .. controls +(0,-0.7) and +(0,0.3) .. (cancelled.west);
  \draw[cancel] (confirmed.south) -- node[right, font=\tiny, color=dangerred]{5 min} (cancelled.north);
\end{tikzpicture}
\end{center}
\vspace{0.3cm}

O cancelamento repõe automaticamente o \textbf{stock} de todos os produtos. Para encomendas no estado \texttt{confirmed}, o cancelamento só é permitido nos primeiros \textbf{5 minutos}.

\subsection{Notificações por Email}

São enviadas notificações automáticas via Nodemailer (Gmail SMTP) nos seguintes eventos:

\begin{center}
\begin{tabular}{lll}
\toprule
\textbf{Evento} & \textbf{Destinatário} & \textbf{Trigger} \\
\midrule
Mudança de estado de encomenda & Cliente & Qualquer transição de estado \\
Aprovação / rejeição de registo & Supermercado & Ação do admin \\
Nova entrega disponível & Todos os estafetas livres & Estado \texttt{awaiting\_courier} \\
\bottomrule
\end{tabular}
\end{center}

\subsection{Funcionalidades Extra (Bonificação)}

\begin{infobox}
As funcionalidades abaixo constituem extras implementados para além do enunciado base, elegíveis para bonificação.
\end{infobox}

\vspace{0.3cm}
\textbf{Mapa Interativo}

A página \texttt{/admin/map} apresenta todos os supermercados aprovados georreferenciados, utilizando \textbf{Leaflet.js} com tiles OpenStreetMap. As coordenadas são obtidas automaticamente via \textbf{Nominatim} (geocodificação reversa) no momento do registo ou edição da morada, sem intervenção manual.

\vspace{0.3cm}
\textbf{Cupões de Desconto}

Sistema completo de cupões com as seguintes características:
\begin{itemize}[leftmargin=*, label=\textcolor{primaryblue}{\textbullet}]
  \item Criados pelo \textbf{supermercado} e válidos apenas na sua loja;
  \item Apenas desconto por percentagem (1\%--100\%);
  \item Controlo de data de expiração e número máximo de usos;
  \item Ativação/desativação pelo supermercado e pelo admin;
  \item O admin pode visualizar e editar todos os cupões da plataforma, com informação do supermercado criador;
  \item Aplicados em vendas presenciais (\texttt{/supermarket/newsale}) com desconto calculado em tempo real.
\end{itemize}

\vspace{0.3cm}
\textbf{Notificações por Email}

Descrito na secção anterior --- envio automático em todos os eventos relevantes do ciclo de vida de encomendas e supermercados.

\vspace{0.3cm}
\textbf{API Documentada com Swagger}

Documentação interativa completa em \texttt{/api-docs} com OpenAPI 3.0: todos os endpoints com esquemas de request/response, exemplos e autenticação via cookie JWT.

% ══════════════════════════════════════════════════════════════════════════════
\section{Como Colocar a Solução a Executar}

\subsection{Pré-requisitos}

\begin{itemize}[leftmargin=*, label=\textcolor{primaryblue}{\textbullet}]
  \item Node.js $\geq$ 18;
  \item MongoDB local ou MongoDB Atlas;
  \item Conta Gmail com \textit{App Password} para envio de emails (opcional).
\end{itemize}

\subsection{Instalação e Execução}

\begin{enumerate}
  \item Clonar o repositório:
\begin{lstlisting}
git clone <url-do-repositorio>
cd paw-group-13
\end{lstlisting}

  \item Instalar as dependências:
\begin{lstlisting}
npm install
\end{lstlisting}

  \item Criar o ficheiro \texttt{.env} na raiz do projeto:
\begin{lstlisting}
MONGO_URI=mongodb+srv://<user>:<pass>@<cluster>.mongodb.net/<db>
secret=<chave-secreta-jwt>
EMAIL_USER=<email@gmail.com>
EMAIL_PASS=<app-password-gmail>
\end{lstlisting}

  \item Iniciar o servidor em modo de desenvolvimento:
\begin{lstlisting}
npm start
\end{lstlisting}

  \item Aceder à aplicação em \url{http://localhost:3000}. A documentação Swagger está em \url{http://localhost:3000/api-docs}.
\end{enumerate}

\subsection{Dependências Principais}

\begin{center}
\begin{tabular}{lll}
\toprule
\textbf{Pacote} & \textbf{Versão} & \textbf{Utilização} \\
\midrule
express               & \textasciitilde4.16 & Framework web \\
mongoose              & \textasciicircum9.3 & ODM MongoDB \\
jsonwebtoken          & \textasciicircum9.0 & Autenticação JWT \\
bcrypt                & \textasciicircum6.0 & Hash de passwords \\
ejs                   & \textasciitilde2.6  & Motor de templates \\
multer                & \textasciicircum2.1 & Upload de imagens \\
nodemailer            & \textasciicircum8.0 & Envio de emails \\
express-mongo-sanitize & \textasciicircum2.2 & Segurança NoSQL \\
swagger-ui-express    & \textasciicircum5.0 & Documentação API \\
swagger-jsdoc         & \textasciicircum6.2 & Geração OpenAPI \\
dotenv                & \textasciicircum17  & Variáveis de ambiente \\
\bottomrule
\end{tabular}
\end{center}

% ══════════════════════════════════════════════════════════════════════════════
\section{Logins de Teste}

Os seguintes utilizadores estão disponíveis na base de dados de teste:

\begin{center}
\begin{tabular}{llll}
\toprule
\textbf{Perfil} & \textbf{Email} & \textbf{Password} & \textbf{Notas} \\
\midrule
Superadmin   & \texttt{superadmin@gmail.com}      & \texttt{estg2026}   & Criado automaticamente; não editável nem eliminável \\
Admin        & \texttt{admin@marketplace.pt}      & \texttt{admin123}   & Acesso total \\
Supermercado & \texttt{mercado@marketplace.pt}    & \texttt{super123}   & Estado: aprovado \\
Estafeta     & \texttt{estafeta@marketplace.pt}   & \texttt{courier123} & Disponível \\
Cliente      & \texttt{cliente@marketplace.pt}    & \texttt{cliente123} & Conta ativa \\
\bottomrule
\end{tabular}
\end{center}

\begin{infobox}
\textbf{Nota:} A conta \texttt{superadmin@gmail.com} é criada automaticamente na primeira inicialização do servidor — não é necessário inserção manual. As restantes credenciais devem ser substituídas pelas contas reais presentes na base de dados de teste antes da entrega.
\end{infobox}

\subsection{Percursos de Teste Recomendados}

\subsubsection{Admin}
\begin{enumerate}[leftmargin=*, label=\arabic*.]
  \item Login em \texttt{/auth/login} com conta admin;
  \item Verificar o dashboard (\texttt{/admin/dashboard}) com contagens globais;
  \item Aprovar um supermercado pendente em \texttt{/admin/approvals} --- verificar email de aprovação;
  \item Criar um estafeta e um cliente em \texttt{/admin/users};
  \item Criar uma categoria de produto em \texttt{/admin/categories};
  \item Consultar os cupões da plataforma em \texttt{/admin/coupons} e ver o detalhe de um;
  \item Visualizar o mapa de supermercados em \texttt{/admin/map};
  \item Testar a API em \texttt{/api-docs} (fazer login primeiro para o cookie estar ativo).
\end{enumerate}

\subsubsection{Supermercado}
\begin{enumerate}[leftmargin=*, label=\arabic*.]
  \item Login com conta de supermercado aprovado;
  \item Configurar a loja em \texttt{/supermarket/settings} (morada, horário, métodos de entrega);
  \item Adicionar um produto em \texttt{/supermarket/products} com imagem e categoria;
  \item Registar uma venda presencial em \texttt{/supermarket/newsale}: pesquisar cliente por email, adicionar produtos, aplicar cupão (opcional), finalizar;
  \item Gerir o estado da encomenda criada em \texttt{/supermarket/orders};
  \item Criar um cupão da loja em \texttt{/supermarket/coupons}.
\end{enumerate}

\subsubsection{Estafeta}
\begin{enumerate}[leftmargin=*, label=\arabic*.]
  \item Login com conta de estafeta;
  \item Verificar entregas disponíveis em \texttt{/courier/deliveries} (requer encomenda no estado \texttt{awaiting\_courier});
  \item Aceitar a entrega e confirmá-la em \texttt{/courier/order/:id};
  \item Consultar o histórico de entregas em \texttt{/courier/history}.
\end{enumerate}

% ══════════════════════════════════════════════════════════════════════════════
\section{Considerações Finais}

Na Milestone 1 foram implementadas as funcionalidades base da plataforma, com foco no \textbf{backoffice}: gestão interna pelos perfis \textit{admin}, \textit{supermarket} e \textit{courier}. A API REST está completa e documentada com Swagger, o sistema de autenticação com JWT em cookie garante separação de acessos entre perfis, e os extras de bonificação (mapa georreferenciado, cupões de desconto e notificações por email) estão totalmente implementados.

Para a \textbf{Milestone 2} está previsto o desenvolvimento do frontoffice para o cliente final: navegação no catálogo de supermercados, encomenda online com carrinho e cupões, e histórico de encomendas com cancelamento.

\end{document}
